% !TeX root = ../tcc.tex
\chapter{Referencial teórico}
\label{cap:referencial-teorico}

Neste capítulo são apresentados os conceitos que fundamentam a investigação.
Inicialmente, são abordadas a pesquisa empírica em engenharia de software
e a revisão sistemática da literatura. Em seguida, são apresentados aspectos
da hanseníase, da aplicação de aprendizado de máquina a imagens de pele e
das medidas usadas na avaliação dos modelos.

\section{Pesquisa empírica em engenharia de software}

A pesquisa empírica utiliza observações e dados para investigar fenômenos,
avaliar propostas e fundamentar conclusões. Na engenharia de software, as
estratégias empíricas incluem experimentos, estudos de caso e estudos
secundários de literatura, cada qual com objetivos e procedimentos próprios
\cite{wohlin2024}. A escolha entre essas estratégias depende da questão
investigada e das evidências necessárias para respondê-la.

Um experimento controlado envolve a definição de fatores, tratamentos e
medidas para investigar seus efeitos sob condições estabelecidas. Já uma
revisão sistemática é um estudo secundário, no qual as evidências de outros
trabalhos são identificadas e analisadas a partir de um processo documentado
\cite[p. 51--56]{wohlin2024}. Neste TCC, a revisão busca reunir informações
sobre avaliações de modelos; sua realização não constitui, por si só, um
experimento controlado ou uma validação clínica.

\subsection{Revisão sistemática da literatura}

Uma revisão sistemática parte de questões de pesquisa e de um protocolo
que orienta a busca, a seleção, a avaliação e a síntese dos estudos. Entre
os elementos a registrar estão as fontes e estratégias de busca, os critérios
de seleção e os procedimentos de extração de dados
\cite{kitchenham2007,wohlin2024}. A descrição dessas decisões permite
examinar o alcance e as limitações do conjunto de estudos identificado.

No presente trabalho, a revisão é empregada para compreender como modelos
aplicados a imagens de pele relacionadas à hanseníase são avaliados. O
protocolo e o andamento de sua execução são apresentados no
Capítulo~\ref{cap:revisao-literatura}.

\section{Hanseníase e identificação de manifestações cutâneas}

A hanseníase pode apresentar manifestações na pele e comprometimento de
nervos periféricos. Segundo a \citeonline{oms2018hanseniase}, entre os
sinais usados para o diagnóstico estão a perda de sensibilidade em uma
mancha cutânea, alterações em nervos periféricos associadas à perda de
sensibilidade e a presença de bacilos em exame de esfregaço cutâneo.
Portanto, uma fotografia de pele representa somente parte das informações
consideradas na avaliação clínica.

No contexto desta pesquisa, a análise computacional de imagens é tratada
como possível apoio à triagem. Seu desempenho em uma base de imagens não
demonstra, isoladamente, que o sistema possa estabelecer um diagnóstico ou
melhorar o atendimento em um serviço de saúde. Por isso, é necessário
relacionar os resultados reportados à tarefa realizada, aos dados utilizados
e à forma de validação \cite{fernandes2024}.

\section{Aprendizado de máquina e visão computacional}

A visão computacional compreende métodos capazes de extrair e analisar
informações de imagens. No aprendizado de máquina, modelos são ajustados a
partir de dados para executar uma tarefa; no aprendizado profundo, redes com
múltiplas camadas aprendem representações em diferentes níveis. Redes
neurais convolucionais são uma das abordagens empregadas na análise de
imagens \cite{lecun2015deep}.

As aplicações à hanseníase diferem quanto à tarefa e ao resultado produzido.
Na classificação, o modelo atribui uma categoria à imagem; na detecção,
identifica e localiza objetos ou lesões; na segmentação, associa regiões ou
pixels a classes de interesse. Essas saídas respondem a perguntas diferentes
e devem ser avaliadas com medidas apropriadas à tarefa
\cite{yotsu2023,eita2025,sanchez2024}.

Por exemplo, \citeonline{yotsu2023} investigaram a classificação de doenças
tropicais com manifestações cutâneas, enquanto \citeonline{sanchez2024}
estudaram a segmentação de feridas em imagens de pessoas com hanseníase.
Delimitar uma ferida e identificar a doença em uma imagem são objetivos
distintos; assim, um resultado de segmentação não deve ser apresentado como
medida de diagnóstico da hanseníase.

\section{Métricas e dimensões de avaliação}\label{muxe9tricas-e-dimensuxf5es-de-avaliauxe7uxe3o-do-gqm}

As métricas serão extraídas conforme a tarefa e o protocolo de cada
estudo. Não se exige que todos os artigos apresentem todas as medidas.
Métricas ausentes serão registradas como não reportadas; cálculos
adicionais somente serão realizados quando houver dados suficientes, com
identificação explícita de que foram derivados nesta revisão.

\subsection{Resultados de
classificação}\label{resultados-de-classificauxe7uxe3o}

Para classificação binária, deve-se definir previamente o que constitui
a classe positiva e o diagnóstico/rótulo de referência. Em modelos
multiclasse, a hanseníase pode ser analisada como classe positiva frente
às demais classes (um contra todos). As definições abaixo se referem à
unidade avaliada --- por exemplo, imagem --- e não implicam
automaticamente diagnóstico por paciente.

{\def\LTcaptype{none} % do not increment counter
\begin{longtable}[]{@{}
  >{\raggedright\arraybackslash}p{(\linewidth - 2\tabcolsep) * \real{0.5000}}
  >{\raggedright\arraybackslash}p{(\linewidth - 2\tabcolsep) * \real{0.5000}}@{}}
\toprule\noalign{}
\begin{minipage}[b]{\linewidth}\raggedright
Resultado
\end{minipage} & \begin{minipage}[b]{\linewidth}\raggedright
Significado
\end{minipage} \\
\midrule\noalign{}
\endhead
\bottomrule\noalign{}
\endlastfoot
\textbf{TP --- True Positive} & A unidade pertence à classe positiva
segundo a referência e o modelo a classifica como positiva. \\
\textbf{TN --- True Negative} & A unidade pertence à classe negativa
segundo a referência e o modelo a classifica como negativa. \\
\textbf{FP --- False Positive} & A unidade pertence à classe negativa
segundo a referência, mas o modelo a classifica como positiva. \\
\textbf{FN --- False Negative} & A unidade pertence à classe positiva
segundo a referência, mas o modelo a classifica como negativa. \\
\end{longtable}
}

A classe positiva pode representar hanseníase ou uma manifestação
definida no protocolo do estudo. Essas tarefas devem ser distinguidas:
reconhecer uma ferida em paciente com hanseníase não equivale a
diagnosticar a doença.

\subsection{Métricas relacionadas à questão de
pesquisa}\label{muxe9tricas-relacionadas-uxe0-questuxe3o-de-pesquisa}

{\def\LTcaptype{none} % do not increment counter
\begin{longtable}[]{@{}
  >{\raggedright\arraybackslash}p{(\linewidth - 6\tabcolsep) * \real{0.2500}}
  >{\raggedright\arraybackslash}p{(\linewidth - 6\tabcolsep) * \real{0.2500}}
  >{\raggedright\arraybackslash}p{(\linewidth - 6\tabcolsep) * \real{0.2500}}
  >{\raggedright\arraybackslash}p{(\linewidth - 6\tabcolsep) * \real{0.2500}}@{}}
\toprule\noalign{}
\begin{minipage}[b]{\linewidth}\raggedright
ID
\end{minipage} & \begin{minipage}[b]{\linewidth}\raggedright
Dimensão do desempenho
\end{minipage} & \begin{minipage}[b]{\linewidth}\raggedright
Métrica
\end{minipage} & \begin{minipage}[b]{\linewidth}\raggedright
Descrição e operacionalização
\end{minipage} \\
\midrule\noalign{}
\endhead
\bottomrule\noalign{}
\endlastfoot
M1 & Classificação & \textbf{Sensibilidade --- Sensitivity (SEN/Recall)}
& Proporção de positivos corretamente identificados: \textbf{SEN = TP /
(TP + FN)}. \\
M2 & Classificação & \textbf{Especificidade --- Specificity (SPE)} &
Proporção de negativos corretamente identificados: \textbf{SPE = TN /
(TN + FP)}. \\
M3 & Classificação & \textbf{Precisão --- Precision (PRE)} & Proporção
de predições positivas que correspondem a positivos de referência:
\textbf{PRE = TP / (TP + FP)}. \\
M4 & Classificação & \textbf{F1-score} & Média harmônica entre precisão
e sensibilidade, considerando o efeito de falsos positivos e falsos
negativos: \textbf{F1 = 2 $\times$ PRE $\times$ SEN / (PRE + SEN)}, ou \textbf{2TP /
(2TP + FP + FN)}. \\
M5 & Classificação geral & \textbf{Acurácia --- Accuracy (ACC)} &
Proporção total de classificações corretas: \textbf{ACC = (TP + TN) /
(TP + TN + FP + FN)}, na formulação binária. Para multiclasse, usar
acertos totais / unidades totais. \\
M6 & Equidade do desempenho & \textbf{Diferença de desempenho entre
grupos} & Comparação da mesma métrica entre grupos de tons de pele,
quando disponível. Para dois grupos A e B: \textbf{$\Delta$ = métrica(A) $-$
métrica(B)}, com grupos e direção identificados. Registrar tamanhos dos
grupos, condições de avaliação e incerteza, quando reportados. \\
M7 & Generalização & \textbf{Variação de desempenho entre datasets} &
Comparação da mesma métrica entre avaliações em dados de distribuição
semelhante à de treinamento e dados de distribuição diferente:
\textbf{gap = métrica(ID) $-$ métrica(OOD)}, quando o desenho permitir.
Registrar se houve adaptação ou retreinamento e se tarefa/classes são
comparáveis. \\
M8 & Segmentação / sobreposição espacial & \textbf{Intersection over
Union (IoU)} & Interseção dividida pela união das regiões prevista P e
de referência G: \textbf{IoU = área(P $\cap$ G) / área(P $\cup$ G)}. Em máscaras
binárias: \textbf{IoU = TP\_pixel / (TP\_pixel + FP\_pixel +
FN\_pixel)}. Registrar mIoU separadamente quando houver média entre
classes. \\
M9 & Detecção / segmentação por instância & \textbf{Mean Average
Precision (mAP)} & Média dos valores de AP das classes consideradas; AP
resume a curva precisão--recall segundo a convenção do benchmark.
Registrar o critério de correspondência e o limiar de IoU: por exemplo,
\textbf{mAP@0,50} ou \textbf{mAP@0,50:0,95}. Distinguir AP de caixas e
AP de máscaras. \\
\end{longtable}
}

\textbf{IoU e mAP não são métricas exclusivas de segmentação.} Em
detecção, IoU mede sobreposição de caixas e pode determinar a
correspondência entre predição e referência; mAP resume o desempenho das
detecções. Na segmentação semântica, IoU/mIoU mede sobreposição de
máscaras por classe. Na segmentação por instância, também se utiliza
AP/mAP calculada sobre máscaras. A documentação de avaliação COCO
explicita essa distinção entre caixas e máscaras \cite{coco}.

\textbf{Critérios de interpretação:} registrar se a unidade é paciente,
imagem, objeto ou pixel; separar resultados específicos de hanseníase de
médias globais; identificar médias macro, micro ou ponderadas; registrar
denominadores nulos como indefinidos ou conforme a convenção explicitada
no estudo. Diferenças entre métricas expressas em porcentagem devem ser
apresentadas em pontos percentuais. M6 e M7 são dimensões de comparação,
não métricas únicas universalmente definidas. Desempenhos de modelos
diferentes em bases diferentes não isolam o efeito da mudança de
distribuição.

O conjunto de métricas funciona como orientação para extração, não como
lista fechada. Medidas adicionais reportadas pelos artigos, como MCC,
devem ser registradas com sua definição e contexto. A análise de
equidade não será restrição de elegibilidade dos artigos e não será
presumida quando não houver resultados por grupo.

